% GENERATED FILE. Edit canonical lessons, then run npm run generate:books.
% canonical-source-hash: fnv1a64:a4e932e8ead3d6db
% canonical-lessons: TA-C194-urru, TA-C194-narukku, TA-C194-uri, TA-C194-pori, TA-C194-picai

\chapter{To pour, To chop, To peel, To deep-fry, To knead}
\label{ch:ta-to-pour-to-chop-to-peel-to-deep-fry-to-knead}
\hlchaptermodality{194}

\begin{chapteropening}
\textbf{By the end of this chapter:} \emph{I can say to pour, to chop, to peel, to deep-fry and to knead.}

Complete the last lesson of chapter 194: I can say to pour, to chop, to peel, to deep-fry and to knead.
\end{chapteropening}

\section[\texorpdfstring{\ta{ஊற்று} (ūṟṟu)}{ஊற்று (ūṟṟu)}]{\ta{ஊற்று} (ūṟṟu) — to pour}

\label{lesson:TA-C194-urru}



\begin{quote}
Before the new one: say the Tamil for a coach, a trainer, then the Tamil for a carpenter.
\end{quote}

\subsection*{You'll want to know: \ta{ஊற்று}}
\textbf{\ta{ஊற்று}} — \emph{ūṟṟu} — \textquotedblleft{}to pour\textquotedblright{}.

\begin{grammarlens}[title={What you've built}]
One more word for this chapter.
\end{grammarlens}

\subsection*{Your turn}
\begin{itemize}
\raggedright
  \item \emph{Say it:} \emph{ūṟṟu}
  \item \emph{Say it:} \emph{ūṟṟu}, once more
  \item \emph{Say it:} say \emph{ūṟṟu}
  \item \emph{Recall:} say the Tamil for a coach, a trainer, then the Tamil for a carpenter, then say \emph{ūṟṟu} again
\end{itemize}

\begin{tcolorbox}[breakable,colback=teal!4,colframe=teal!35!black,title={Before you move on}]
What does \ta{ஊற்று} mean? (\textquotedblleft{}To pour\textquotedblright{}.) Say it once more.
\end{tcolorbox}
\section[\texorpdfstring{\ta{நறுக்கு} (naṟukku)}{நறுக்கு (naṟukku)}]{\ta{நறுக்கு} (naṟukku) — to chop}

\label{lesson:TA-C194-narukku}



\begin{quote}
Before the new one: say the Tamil for a carpenter, then the Tamil for to pour.
\end{quote}

\subsection*{You'll want to know: \ta{நறுக்கு}}
\textbf{\ta{நறுக்கு}} — \emph{naṟukku} — \textquotedblleft{}to chop\textquotedblright{}.

\begin{grammarlens}[title={What you've built}]
One more word for this chapter.
\end{grammarlens}

\subsection*{Your turn}
\begin{itemize}
\raggedright
  \item \emph{Say it:} \emph{naṟukku}
  \item \emph{Say it:} \emph{naṟukku}, once more
  \item \emph{Say it:} say \emph{naṟukku}
  \item \emph{Recall:} say the Tamil for a carpenter, then the Tamil for to pour, then say \emph{naṟukku} again
\end{itemize}

\begin{tcolorbox}[breakable,colback=teal!4,colframe=teal!35!black,title={Before you move on}]
What does \ta{நறுக்கு} mean? (\textquotedblleft{}To chop\textquotedblright{}.) Say it once more.
\end{tcolorbox}
\section[\texorpdfstring{\ta{உரி} (uri)}{உரி (uri)}]{\ta{உரி} (uri) — to peel}

\label{lesson:TA-C194-uri}



\begin{quote}
Before the new one: say the Tamil for to pour, then the Tamil for to chop.
\end{quote}

\subsection*{You'll want to know: \ta{உரி}}
\textbf{\ta{உரி}} — \emph{uri} — \textquotedblleft{}to peel\textquotedblright{}.

\begin{grammarlens}[title={What you've built}]
One more word for this chapter.
\end{grammarlens}

\subsection*{Your turn}
\begin{itemize}
\raggedright
  \item \emph{Say it:} \emph{uri}
  \item \emph{Say it:} \emph{uri}, once more
  \item \emph{Say it:} say \emph{uri}
  \item \emph{Recall:} say the Tamil for to pour, then the Tamil for to chop, then say \emph{uri} again
\end{itemize}

\begin{tcolorbox}[breakable,colback=teal!4,colframe=teal!35!black,title={Before you move on}]
What does \ta{உரி} mean? (\textquotedblleft{}To peel\textquotedblright{}.) Say it once more.
\end{tcolorbox}
\section[\texorpdfstring{\ta{பொரி} (pori)}{பொரி (pori)}]{\ta{பொரி} (pori) — to deep-fry}

\label{lesson:TA-C194-pori}



\begin{quote}
Before the new one: say the Tamil for to chop, then the Tamil for to peel.
\end{quote}

\subsection*{You'll want to know: \ta{பொரி}}
\textbf{\ta{பொரி}} — \emph{pori} — \textquotedblleft{}to deep-fry\textquotedblright{}.

\begin{grammarlens}[title={What you've built}]
One more word for this chapter.
\end{grammarlens}

\subsection*{Your turn}
\begin{itemize}
\raggedright
  \item \emph{Say it:} \emph{pori}
  \item \emph{Say it:} \emph{pori}, once more
  \item \emph{Say it:} say \emph{pori}
  \item \emph{Recall:} say the Tamil for to chop, then the Tamil for to peel, then say \emph{pori} again
\end{itemize}

\begin{tcolorbox}[breakable,colback=teal!4,colframe=teal!35!black,title={Before you move on}]
What does \ta{பொரி} mean? (\textquotedblleft{}To deep-fry\textquotedblright{}.) Say it once more.
\end{tcolorbox}
\section[\texorpdfstring{\ta{பிசை} (picai)}{பிசை (picai)}]{\ta{பிசை} (picai) — to knead}

\label{lesson:TA-C194-picai}



\begin{quote}
Before the new one: say the Tamil for to peel, then the Tamil for to deep-fry.
\end{quote}

\subsection*{You'll want to know: \ta{பிசை}}
\textbf{\ta{பிசை}} — \emph{picai} — \textquotedblleft{}to knead\textquotedblright{}.

\begin{grammarlens}[title={What you've built}]
One more word for this chapter.
\end{grammarlens}

\subsection*{Your turn}
\begin{itemize}
\raggedright
  \item \emph{Say it:} \emph{picai}
  \item \emph{Say it:} \emph{picai}, once more
  \item \emph{Say it:} say \emph{picai}
  \item \emph{Recall:} say the Tamil for to peel, then the Tamil for to deep-fry, then say \emph{picai} again
\end{itemize}

\begin{tcolorbox}[breakable,colback=teal!4,colframe=teal!35!black,title={Before you move on}]
What does \ta{பிசை} mean? (\textquotedblleft{}To knead\textquotedblright{}.) Say it once more.
\end{tcolorbox}
